% Generated draft: requires mathematical and source review.
\documentclass[11pt]{article}
\usepackage[margin=25mm]{geometry}
\usepackage{fontspec}
\setmainfont{Latin Modern Roman}
\usepackage{amsmath,amssymb,mathtools,unicode-math}
\setmathfont{Latin Modern Math}
\usepackage{enumitem,needspace,xurl,hyperref}
\setlength{\parindent}{0pt}
\setlength{\parskip}{5pt}
\setlength{\emergencystretch}{3em}
\newcommand{\N}{\mathbb{N}}
\newcommand{\Z}{\mathbb{Z}}
\newcommand{\Q}{\mathbb{Q}}
\newcommand{\R}{\mathbb{R}}
\newcommand{\C}{\mathbb{C}}
\newcommand{\F}{\mathbb{F}}
\newcommand{\A}{\mathbb{A}}
\newcommand{\PP}{\mathbb{P}}
\newcommand{\E}{\mathbb{E}}
\newcommand{\eps}{\varepsilon}
\newcommand{\dd}{\,\mathrm{d}}
\begin{document}
\section*{UnsolvedMath ID: 30006299}
\textbf{Source problem code:} catalogue entry
\subsection{Definitions and mathematical statement}
The supplied question concerns an arbitrary matroid \(M\), but does not specify finiteness, a realization field, the ambient contact manifold, or the meaning of correspondence. The following definitions describe standard possible interpretations; they do not resolve these missing specifications.

For a finite matroid \(M=(E,r)\), its rank function \(r:2^E\to\Z_{\geq 0}\) satisfies \(0\leq r(A)\leq |A|\), \(r(A)\leq r(B)\) when \(A\subseteq B\), and
\[
r(A\cup B)+r(A\cap B)\leq r(A)+r(B)
\qquad(A,B\subseteq E).
\]
For a specified field \(k\), a labeled vector realization of \(M\) is a family \((v_e)_{e\in E}\) in \(k^{r(E)}\) such that
\[
\dim_k\operatorname{span}_k\{v_e:e\in A\}=r(A)
\qquad(A\subseteq E).
\]
Write \(\operatorname{Real}^{\mathrm{vec}}_k(M)\) for the set of these families. Passing to equivalence classes under changes of basis, additionally allowing individual nonzero column rescalings, or retaining a realization scheme gives different possible targets. The source does not select one. It also does not say whether only finite matroids are intended.

A Legendrian submanifold of a contact manifold \((Y^{2n+1},\xi)\) is an \(n\)-dimensional submanifold \(\Lambda\) with \(T\Lambda\subseteq\xi\). In a differential graded algebra interpretation of augmentations, after specifying an associated unital differential graded algebra \((\mathcal A(\Lambda),\partial)\) and its coefficients and grading, an augmentation to a field \(k\) is a unital graded algebra map
\[
\epsilon:\mathcal A(\Lambda)\longrightarrow k,
\qquad \epsilon\circ\partial=0,
\]
where \(k\) has zero differential and is concentrated in degree zero. Any required compatibility with the coefficient ring must also be specified. The source does not identify the augmentation algebra, grading convention, or equivalence relation on augmentations.

\textbf{Question, with unresolved parameters retained.} For every matroid \(M\), does there exist a Legendrian \(\Lambda_M\) whose augmentations correspond naturally to realizations of \(M\)? Once the missing choices are fixed, the proposed relationship has the schematic form
\[
\operatorname{Aug}_k(\Lambda_M)\ \longleftrightarrow\ \operatorname{Real}_k(M).
\]
Neither the quantifier on \(k\), the meaning of the two sides, nor the structure and naturality required of this correspondence is specified. Thus this display is a formulation template, not a fully determined bijection or moduli-space assertion. The supplied edited statement asks for a construction for every matroid, or, alternatively, a precise characterization of the matroids admitting such a model. The original excerpt states the existence question without explicitly giving that alternative.

The original excerpt also says that evidence from symplectic geometry suggests that \(P\) will only be defined after choosing an augmentation for the Legendrian link of the cusp locus of \(L_{\Sigma(1)}\), and that the Donaldson divisor theorem suggests that \(\mathrm{NC}\) neighborhoods may arise via deformations. The supplied record does not define \(P\), \(L_{\Sigma(1)}\), or \(\mathrm{NC}\), or specify how these contextual remarks constrain the question. They are retained as source qualifications, not used as additional hypotheses or asserted conclusions.
\subsection{Short English statement}
The question asks whether every matroid can be modeled by a Legendrian whose augmentations naturally encode its realizations, or which matroids admit such a model. A precise conjecture requires choices of realization field, realization equivalence, contact setting, augmentation theory, and the meaning of natural correspondence. The supplied record labels the problem open; this draft does not independently verify that status.
\subsection{Sources}
\begin{itemize}
\item Local UnsolvedMath registry, permanent ID \(30006299\), title \emph{Legendrian Augmentations and Matroid Realizations}, source problem code \texttt{catalogue entry}. The supplied record belongs to the Oberwolfach Reports Open Problems collection and gives proposed year 2025; no individual proposer is supplied.
\item \emph{Toric Geometry} (2025), Oberwolfach Reports, pp.~883--940, DOI 10.4171/owr/2025/19. Supplied source and literature-source link: \url{https://doi.org/10.4171/owr/2025/19}. These reports document mathematical workshops and their open-problem discussions. No exact page or named contribution for this question is supplied.
\item Supplied literature assessment, checked on \textbf{2026-08-22}: the record labels the problem open and reports that the report poses the correspondence as a general question and that no authoritative source located by that cutoff proves it for arbitrary matroids. This is supplied dated curation derived from the OWR record, not an independently verified literature assessment. The listed source checked is the Oberwolfach report linked above; time-sensitive status and cited claims require independent verification.
\item The registry labels the statement \texttt{corrected\_verified} and states that only clear formula delimiters, notation, grammar, and typography were repaired while preserving mathematical content and scope. This is a supplied editorial qualification, not independent mathematical verification of this draft.
\end{itemize}
\end{document}
